\documentclass[10pt,a4paper]{amsart}
\usepackage[margin=19mm]{geometry}
\usepackage{amsmath,amssymb,mathtools}
\usepackage[scr=rsfs,cal=euler]{mathalfa}
\usepackage{xfrac}
\usepackage[unicode,hidelinks,bookmarks=false]{hyperref}
\newcommand{\ie}{i.e.\ }
\newcommand{\cf}{cf.\ }
\newcommand{\resp}{respectively\ }
\newcommand{\Subsection}{\subsection}
\providecommand{\nobreakdash}{\nobreak\hbox{-}}
\setlength{\emergencystretch}{2em}
\multlinegap0pt
\usepackage{amssymb}
\usepackage{mathtools}
\usepackage{bm}
\usepackage[shortlabels]{enumitem}
\newtheorem{lemma}{Lemma}
\newtheorem{theorem}{Theorem}
\usepackage{cleveref}
\crefname{lemma}{Lemma}{Lemmas}
\crefname{theorem}{Theorem}{Theorems}
\title{Dual Geometry of Spherical Designs:\\ Polarity, Self-Polar Rigidity, and Quadrature Structure}
\author{Congpei An}
\date{Source: arXiv:2609.02439v1 (CC BY 4.0)}
\begin{document}
\maketitle

\noindent\emph{Source and licence.} Dual Geometry of Spherical Designs:\\ Polarity, Self-Polar Rigidity, and Quadrature Structure. Version arXiv:2609.02439v1 (CC BY 4.0), distributed by arXiv under the Creative Commons Attribution 4.0 International licence, http://creativecommons.org/licenses/by/4.0/, as recorded in the arXiv licence metadata for this exact version and shown on the abstract page https://arxiv.org/abs/2609.02439v1. Full licence notice: \texttt{licenses/arxiv-2609.02439-CC-BY-4.0.txt}.\par\smallskip

\noindent\textit{Editorial scope.} Counted unit: the article's theorem thm:stable-quadrature (source0.tex lines 801--829). The article's complete original proof of it is delivered with it (source0.tex lines 831--873, one \texttt{proof} environment of 43 lines). No mathematics is written, repaired or re-derived: the statement and the proof are the article's own lines, in the article's own notation, and no retained line is substituted or re-flowed.  Retained uncounted as the source-local context the proof needs: lines 744--758.  Nothing was omitted from any retained span, so the package carries no omission record.  No retained line contains a citation, so the wrapper carries no bibliography and no bibliography entry.  No retained line contains a figure environment, an image-inclusion command or drawing code, so no image is delivered and the package declares no asset dependency.  Editorial formatting only, in addition: an AMS \texttt{amsart} preamble that restates the article's own declarations and macros used by the retained lines, namely the environments \texttt{lemma}, \texttt{theorem} on separate counters, together with the standard packages those lines need.  The retained environments print in this document as Lemma 1, Theorem 1 in order of appearance, whereas the article prints the same items with the numbering of its own full document.  The complete native source of the article as distributed in the arXiv e-print for this exact version is bundled unchanged as \texttt{sources/209200/source0.tex}.
\begin{lemma}[Coercivity of the extremal certificate]\label{lem:certificate-coercive}
There exist constants $\varepsilon_0,c_0>0$, depending only on $d$ and $t$, with the following property.  For every
\[
 \eta_{t,d}\le\alpha\le\eta_{t,d}+\varepsilon_0,
\]
there is a node set $\Xi_\alpha$ obtained from $\Xi_{t,d}$ by replacing its largest node $\eta_{t,d}$ with $\alpha$, and a nonnegative polynomial $\Phi_\alpha$ on $[-1,\alpha]$ such that
\begin{equation}\label{eq:certificate-coercive}
 \Phi_\alpha(s)\ge c_0\,\operatorname{dist}(s,\Xi_\alpha)^2,
 \qquad -1\le s\le\alpha.
\end{equation}
Moreover, if $\delta:=\alpha-\eta_{t,d}$ and $\nu$ is any probability measure supported on $[-1,\alpha]$ whose moments agree with $\mu_d$ through degree $t$, then
\begin{equation}\label{eq:certificate-defect}
 \int\Phi_\alpha\,d\nu=\delta.
\end{equation}
\end{lemma}

\begin{theorem}[Stable extremal quadrature]\label{thm:stable-quadrature}
Fix $d\ge2$ and $t\ge2$.  There exist constants $\varepsilon_0,C>0$, depending only on $d$ and $t$, such that the following holds.  Let $\nu$ be a probability measure supported on $[-1,\alpha]$, where
\[
 \eta_{t,d}\le\alpha\le\eta_{t,d}+\varepsilon_0,
 \qquad
 \delta:=\alpha-\eta_{t,d},
\]
and suppose that $\nu$ and $\mu_d$ have the same moments through degree $t$.  Then
\begin{equation}\label{eq:mean-square-clustering}
 \boxed{
 \int \operatorname{dist}(s,\Xi_{t,d})^2\,d\nu(s)
 \le C\delta,}
\end{equation}
and
\begin{equation}\label{eq:W1-stable}
 \boxed{
 W_1(\nu,\nu_{t,d}^*)\le C\sqrt\delta.}
\end{equation}

More precisely, choose the nearest-node projection from $[-1,\alpha]$ onto the moving set $\Xi_\alpha$ of \cref{lem:certificate-coercive}, and let $E_k(\alpha)$ be its Voronoi cells.  If
\[
 w_k:=\nu(E_k(\alpha)),
\]
with the top cell indexed so that it is mapped from $\alpha$ back to $\eta_{t,d}$, then
\begin{equation}\label{eq:weight-stability}
 \boxed{
 \sum_{k=0}^n|w_k-\lambda_k|\le C\sqrt\delta.}
\end{equation}
\end{theorem}

\begin{proof}
By \cref{lem:certificate-coercive},
\[
 c_0\int\operatorname{dist}(s,\Xi_\alpha)^2\,d\nu(s)
 \le\int\Phi_\alpha\,d\nu
 =\delta.
\]
Let $\pi_\alpha$ be a measurable nearest-node projection and set
\[
 \widetilde\nu:=(\pi_\alpha)_\#\nu.
\]
The coupling $s\mapsto\pi_\alpha(s)$ gives
\begin{equation}\label{eq:projection-W1}
 W_1(\nu,\widetilde\nu)
 \le\left(\int|s-\pi_\alpha(s)|^2\,d\nu(s)\right)^{1/2}
 \le C\sqrt\delta.
\end{equation}
Move the top atom of $\widetilde\nu$ from $\alpha$ to $\eta_{t,d}$ and leave the remaining atoms fixed.  Denote the resulting measure by
\[
 \widehat\nu=\sum_{k=0}^n w_k\delta_{\xi_k}.
\]
Then
\[
 W_1(\widetilde\nu,\widehat\nu)\le\delta,
\]
so
\begin{equation}\label{eq:nu-hat-W1}
 W_1(\nu,\widehat\nu)\le C\sqrt\delta.
\end{equation}

It remains to compare the weights $w_k$ with the quadrature weights.  Let
\[
 V=(\xi_k^\ell)_{0\le\ell,k\le n}
\]
be the Vandermonde matrix of the distinct quadrature nodes.  Since $V$ is invertible, it suffices to compare the first $n$ moments.  For $1\le\ell\le n$, the function $s\mapsto s^\ell$ is $\ell$-Lipschitz on $[-1,1]$.  Hence, using moment matching of $\nu$ and exactness of $\nu_{t,d}^*$,
\[
 \left|\sum_k(w_k-\lambda_k)\xi_k^\ell\right|
 =\left|\int s^\ell\,d(\widehat\nu-\nu)\right|
 \le \ell W_1(\widehat\nu,\nu)
 \le C\sqrt\delta.
\]
The zeroth moment difference is zero.  Applying $V^{-1}$ proves \eqref{eq:weight-stability}.  Since two probability measures on the same finite node set have Wasserstein distance bounded by their $\ell^1$ weight difference, \eqref{eq:W1-stable} follows from \eqref{eq:nu-hat-W1} and \eqref{eq:weight-stability}.  Finally, $\operatorname{dist}(s,\Xi_{t,d})\le \operatorname{dist}(s,\Xi_\alpha)+\delta$.  Hence, after squaring and integrating, the bound for $\Xi_\alpha$ gives \eqref{eq:mean-square-clustering} (shrinking $\varepsilon_0$ so that $\delta\le1$).  \qedhere
\end{proof}

\end{document}
